\par\Needspace{10\baselineskip}\pdfbookmark[1]{Pasquich. 1817 Oktober 4.}{pasquich1817}
\begin{center}
\rule{0.65\linewidth}{0.4pt}\\
Göttingische gelehrte Anzeigen. 1817 Oktober 4.\\
\rule{0.65\linewidth}{0.4pt}
\end{center}

\textit{Abgekürzte Logarithmisch-Trigonometrische Tafeln, mit neuen Zusätzen zur
Abkürzung und Erleichterung trigonometrischer Rechnungen, herausgegeben von }\textsc{Joh.}
\textsc{Pasquich}\textit{, Director der Königl. Ofner Sternwarte.} Leipzig. In der Weidmannschen Buchhandlung 1817. XXXII und 228 Seiten in Oktav. (Auch mit lateinischem Titel.)
\par\SourceLeaf{256}
Kleinere Logarithmentafeln mit fünf Dezimalstellen sind bei allen, die viel mit Zahlen rechnen, besonders bei Astronomen, sehr beliebt. Denn die Fälle, in denen sie ausreichen, sind tatsächlich häufig, ja die häufigeren; ein bequemes Format und ein mäßiger Umfang erleichtern die Arbeit sehr. Die kleinen sauberen Tafeln von \textsc{Lalande} und die von \textsc{Prasse} sind allgemein verbreitet. Bei Letzteren ist die Unterbringung auf engem Raum so weit wie möglich getrieben, zum Teil allerdings auf Kosten der Bequemlichkeit. Wie die Vorrede berichtet, wurde die Herausgabe der vorliegenden Tafeln, die ebenfalls nur fünf Stellen bieten, durch den Wunsch veranlasst, den \textsc{Gauss} 1812 in der Monatlichen Correspondenz geäußert hatte: Die dort erstmals gedruckte Tafel zur bequemen Berechnung der Logarithmen von Summen und Differenzen möge in eine neue Ausgabe von \textsc{Lalande}s Tafeln aufgenommen werden. Der neue Druck dieser Hilfstafel in der vorliegenden Sammlung wird denen willkommen sein, die den ersten Druck nicht zur Verfügung hatten oder denen ihre Benutzung in der M. C. zu umständlich war. Außerdem zeichnet sich diese Sammlung durch eine neue Hilfstafel aus, die Herr \textsc{Pasquich} berechnet und den trigonometrischen Tafeln beigefügt hat; darauf werden wir unten näher eingehen.

Die Logarithmen der Zahlen reichen wie bei \textsc{Lalande} und von \textsc{Prasse} bis 10000 und sind hier wie dort in ihrer natürlichen Reihenfolge gedruckt. Allerdings vermisst man zwei Erleichterungen, die \textsc{Lalande}s Tafeln bieten: Zum einen fehlen die \textit{Differenzen}, zum anderen werden die untersten Logarithmen jeder Spalte oben in der folgenden nicht wiederholt. Bei \textsc{Prasse}s Tafeln fehlen diese Annehmlichkeiten zwar ebenfalls, doch gleicht der geringe Platzbedarf der Tafel dies mehr als aus: Dort bieten 24 Seiten genau dasselbe, was bei \textsc{Pasquich} auf 56 Seiten in einem erheblich größeren Format steht. Dies scheinen nur Kleinigkeiten zu sein, und für alle, die nur gelegentlich Logarithmen nachschlagen, sind sie es auch; nicht aber für diejenigen, die ständig Logarithmentafeln zur Hand haben müssen.

\Needspace{14\baselineskip}Bei den trigonometrischen Tafeln enthält jede linke Seite die Logarithmen der Sinus, Kosinus, Tangenten und Kotangenten, und zwar so, dass drei Seiten jeweils zwei Grade umfassen. Diese durch das gewählte Format und die Schrift bedingte Anordnung scheint uns etwas unbequem; \SourceLeaf{257}
wir hätten entweder ein kleineres Format mit jeweils einem halben Grad auf einer Seite oder ein wenig längeres mit kleinerer Schrift vorgezogen, sodass ein ganzer Grad auf die Seite gekommen wäre, wie in \textsc{Shervins} Tafeln. Von diesen Logarithmen sind jeweils nur die vier, drei oder zwei letzten Ziffern gedruckt, solange die vorhergehenden unverändert bleiben. Dadurch wurde die Arbeit für den Abschreiber, den Setzer und den Korrektor erleichtert, und die Tafeln erhalten ein saubereres Aussehen. Dennoch können wir diese Anordnung bei Tafeln für den täglichen Gebrauch nicht uneingeschränkt billigen: Das Auge muss sich immer die, wenn auch geringe, Mühe machen, in der Spalte erst nach oben zu gehen und die übrigen Ziffern zu suchen. Die Differenzen der Logarithmen sind hier bereits durch 60 geteilt. Diese Anordnung wurde auch in einigen anderen Tafeln gewählt, um die Interpolation zu erleichtern. Ob dies tatsächlich eine Erleichterung ist, hängt von den Gewohnheiten des Rechners ab. Der Rezensent findet \textsc{Lalande}s Tafeln, in denen die ganzen Differenzen angegeben sind, in dieser Hinsicht zumindest nicht unbequemer. Da die Zahlen, mit denen gerechnet wird, klein sind, führt ein einigermaßen geübter Rechner die zur Interpolation nötigen Operationen leicht im Kopf aus und findet fast immer Gelegenheit, einen Vorteil des jeweiligen Falls zu nutzen. Dabei hat man außerdem die angenehme Gewissheit, das Interpolationsergebnis so genau wie möglich zu erhalten. Bei der von Herrn \textsc{Pasquich} gewählten Anordnung ist dagegen, allgemein gesprochen, der Interpolationsfehler etwas größer; dies hier näher auszuführen, wäre nicht am Platz.

\Needspace{24\baselineskip}Die rechte Seite enthält bei den trigonometrischen Tafeln die Quadrate der Sinus, Kosinus, Tangenten und Kotangenten. Sie sollen die Interpolation erleichtern, wenn aus dem Logarithmus eines Sinus, Kosinus, einer Tangente oder Kotangente der Logarithmus einer der drei anderen trigonometrischen Funktionen bestimmt werden soll, ohne den Bogen selbst zu benötigen. Diese Operation kommt tatsächlich äußerst häufig vor. Beim gewöhnlichen Verfahren erfordert die Interpolation eine Multiplikation und eine Division, während mit der Hilfstafel von Herrn \textsc{Pasquich} eine Multiplikation ausreicht. Für die Interpolation hinreichend genau gilt nämlich:
\SourceLeaf{258}
\[
\begin{aligned}
\Delta\log\cos\varphi&=-\operatorname{tang}\varphi^{2}.\Delta\log\sin\varphi\\
\Delta\log\operatorname{tang}\varphi&=-\Delta\log\operatorname{cotang}\varphi=(1+\operatorname{tang}\varphi^{2}).\Delta\log\sin\varphi\\
\Delta\log\sin\varphi&=-\operatorname{cotang}\varphi^{2}.\Delta\log\cos\varphi\\
\Delta\log\operatorname{tang}\varphi&=-\Delta\log\operatorname{cotang}\varphi=-(1+\operatorname{cotang}\varphi^{2}).\Delta\log\cos\varphi\\
\Delta\log\sin\varphi&=\cos\varphi^{2}.\Delta\log\operatorname{tang}\varphi=-\cos\varphi^{2}.\Delta\log\operatorname{cotang}\varphi\\
\Delta\log\cos\varphi&=-\sin\varphi^{2}.\Delta\log\operatorname{tang}\varphi=\sin\varphi^{2}.\Delta\log\operatorname{cotang}\varphi
\end{aligned}
\]
Der Rezensent muss jedoch gestehen, dass er das gewöhnliche Interpolationsverfahren dennoch nicht nur ebenso bequem, sondern sogar bequemer findet. Zum einen kostet es stets einige Mühe, sich die sechs obigen Formeln so zur Gewohnheit zu machen, dass man sie ohne jedes Nachdenken oder ohne ein besonderes Blatt daneben richtig anwendet. Zum anderen ist es umständlich, den Multiplikationsfaktor erst auf der anderen Seite nachzuschlagen oder vielmehr zusammenzusuchen, da auch hier die oben erwähnte Trennung der ersten und letzten Ziffern beim Druck gewählt wurde. Schließlich hat man beim gewöhnlichen Verfahren immer nur mit kleinen Zahlen zu tun, mit denen man leicht im Kopf rechnet. Die Quadrate in \textsc{Pasquich}s Tafeln sind dagegen mit fünf Dezimalstellen angegeben; natürlich braucht man diese nicht alle, doch gerade deshalb stören sie, wie jeder erfahrene Rechner weiß. Außerdem können wir nicht unerwähnt lassen, dass das gewöhnliche Verfahren, allgemein gesprochen, \textit{genauer} ist als diese aufwendigere Interpolation. Die Gründe für diese Behauptung, bei der man nach einer weniger genauen Prüfung vielleicht gerade das Gegenteil annehmen könnte, würden hier zu viel Raum beanspruchen. Wir begnügen uns damit, das Gesagte durch ein Beispiel zu erläutern. Soll zu \(\log\cos\varphi=9{,}92478\) der \(\log\operatorname{tang}\varphi\) bestimmt werden, so erhält man den Proportionalteil aus \textsc{Pasquich}s Tafel durch die Berechnung von \(4\times(1+2{,}4170)=13{,}668\) oder gerundet \(=14\), also \(\log\operatorname{tang}\varphi=9{,}80850\). Das gewöhnliche Verfahren liefert den Proportionalteil ebenso bequem durch die Berechnung von \(\frac{4\times28}{9}=12\frac49\), gerundet \(=12\), und den gesuchten Logarithmus \(=9{,}80848\). In diesem Beispiel ist auch das Ergebnis des gewöhnlichen Verfahrens genauer. In anderen Fällen kann das Verhältnis auch umgekehrt sein, doch im Durchschnitt wird der Vorteil in dieser Hinsicht beim gewöhnlichen Verfahren liegen. Wir wollen im Übrigen nicht bestreiten, dass dieser Teil der Tafel, auch wenn die Interpolation dadurch nicht gewinnt, gelegentlich für \SourceLeaf{259}
andere Zwecke willkommen sein kann. Doch die Bequemlichkeit logarithmischer Handtafeln für den täglichen Gebrauch nimmt natürlich in demselben Verhältnis ab, in dem ihr Umfang wächst. Wir bemerken noch, dass im ersten Grad die trigonometrischen Logarithmen von \(10\) zu \(10\) Sekunden bis \(56\) Minuten und in den vier letzten Minuten von \(20\) zu \(20\) Sekunden angegeben sind.

Die Tafel von \textsc{Gauss} für die Logarithmen der Summen und Differenzen ist völlig unverändert abgedruckt. Es scheint uns jedoch inkonsequent, wenn der Verfasser in der Einleitung den Nutzen einer ähnlichen Tafel mit sieben Dezimalstellen bezweifelt. Ist eine solche Tafel zweckmäßig eingerichtet, so ist ihr Nutzen bei genauen Rechnungen genauso groß wie der Nutzen der hier wieder abgedruckten Tafeln bei Rechnungen mit fünf Dezimalstellen. Bei den kleineren wie bei den größeren Tafeln wird die dadurch mögliche Zeitersparnis natürlich nur für diejenigen spürbar, die \textit{viel} zu rechnen haben. Wir dürfen nun bald das Erscheinen einer solchen größeren Tafel erwarten, die von kundiger Hand berechnet ist.

